
```latex
\documentclass[12pt,a4paper]{article}

\usepackage[margin=1in]{geometry}
\usepackage{array}
\usepackage{booktabs}
\usepackage{longtable}
\usepackage{listings}
\usepackage{hyperref}
\usepackage{graphicx}
\usepackage{setspace}

\onehalfspacing

\title{\textbf{Cryptography and Network Security Technical Report}}
\author{Fabiola}
\date{September 2026}

\begin{document}

\maketitle

\tableofcontents
\newpage

\section{Introduction}

This project demonstrates practical applications of cryptography, data
integrity verification, risk assessment and network access control in an
authorised laboratory environment.

The main security objectives of the project are to protect sensitive
information, detect unauthorised modification of data and control access
to network services.

The project includes an encryption and decryption programme, an integrity
verification function using a cryptographic hash, a risk assessment and
a firewall access-control configuration.

\section{Risk Assessment}

\subsection{Assets, Vulnerabilities and Consequences}

The following assets and vulnerabilities were identified during the risk
assessment.

\begin{longtable}{p{0.25\textwidth} p{0.30\textwidth} p{0.35\textwidth}}
\toprule
\textbf{Asset} & \textbf{Vulnerability} & \textbf{Possible Consequence} \\
\midrule
Student records &
Weak staff passwords &
An attacker could gain unauthorised access to student information. \\

Central records server &
Outdated software &
An attacker could exploit known weaknesses and compromise the server. \\

Files transferred between campuses &
Unencrypted file transfers &
Sensitive information could be intercepted during transmission. \\
\bottomrule
\end{longtable}

\subsection{Risk Ranking}

Three major risks were identified and ranked according to likelihood and
impact.

\begin{longtable}{p{0.27\textwidth} p{0.15\textwidth} p{0.15\textwidth} p{0.35\textwidth}}
\toprule
\textbf{Risk} & \textbf{Likelihood} & \textbf{Impact} & \textbf{Reason} \\
\midrule
Unauthorised access to student records &
High &
High &
Weak staff passwords increase the possibility of account compromise,
and access to student records could expose sensitive information. \\

Exploitation of outdated software &
Medium &
High &
Outdated software may contain exploitable weaknesses that could allow
an attacker to compromise the central server. \\

Interception of transferred files &
Medium &
High &
Unencrypted transfers can expose information while files are being
transferred between campuses. \\
\bottomrule
\end{longtable}

\subsection{Recommended Controls}

The following controls were recommended to reduce the identified risks.

\begin{longtable}{p{0.35\textwidth} p{0.50\textwidth}}
\toprule
\textbf{Risk} & \textbf{Recommended Control} \\
\midrule
Unauthorised access to student records &
Enforce strong staff passwords and multi-factor authentication. \\

Exploitation of outdated software &
Regularly update and patch the server software. \\

Interception of transferred files &
Use encrypted file transfers between the campuses. \\
\bottomrule
\end{longtable}

\section{Encryption and Decryption}

Encryption is used to protect information by converting readable
plaintext into ciphertext. The ciphertext should not reveal the original
information to an unauthorised person.

The general process used in the project is:

\begin{center}
\textbf{Plaintext}
$\rightarrow$
\textbf{Encryption}
$\rightarrow$
\textbf{Ciphertext}
\end{center}

Decryption reverses this process and converts the ciphertext back into
the original plaintext.

\begin{center}
\textbf{Ciphertext}
$\rightarrow$
\textbf{Decryption}
$\rightarrow$
\textbf{Original Plaintext}
\end{center}

The programme was tested to verify that information encrypted by the
programme could subsequently be decrypted correctly. A successful
decryption should reproduce the original plaintext.

\subsection{Encryption Test Evidence}

The encryption test follows these steps:

\begin{enumerate}
    \item Input plaintext into the programme.
    \item Encrypt the plaintext.
    \item Observe the resulting ciphertext.
    \item Decrypt the ciphertext.
    \item Compare the decrypted text with the original plaintext.
\end{enumerate}

The expected result is that the decrypted text matches the original
plaintext.

\section{Integrity Verification}

Integrity verification determines whether information has been changed
after it was originally created or transmitted.

The project uses a cryptographic hash function for integrity checking.
A hash converts data into a fixed-length hash value.

The process can be represented as:

\begin{center}
\textbf{Original Data}
$\rightarrow$
\textbf{Hash Function}
$\rightarrow$
\textbf{Hash Value}
\end{center}

If the original data is modified, calculating the hash again should
produce a different hash value.

\subsection{Integrity Test}

The integrity test consists of the following procedure:

\begin{enumerate}
    \item Calculate the hash of the original data.
    \item Store or record the resulting hash value.
    \item Modify the data.
    \item Calculate the hash again.
    \item Compare the two hash values.
\end{enumerate}

If the hash values differ, the data has been modified.

The test therefore demonstrates that cryptographic hashing can be used
to detect changes to information.

\section{Firewall Configuration}

A firewall was designed for an authorised laboratory environment to
demonstrate network access control.

The laboratory network configuration is shown below.

\begin{table}[h]
\centering
\begin{tabular}{ll}
\toprule
\textbf{Component} & \textbf{Configuration} \\
\midrule
Student Records Server & 192.168.10.10 \\
Guest Network & 192.168.20.0/24 \\
Authorised Staff Network & 192.168.30.0/24 \\
Protected Service & SSH \\
Protocol & TCP \\
Port & 22 \\
\bottomrule
\end{tabular}
\caption{Laboratory firewall environment}
\end{table}

The firewall policy consists of three main access-control requirements.

\subsection{Guest Network Blocking}

Access from the guest network to the Student Records Server is blocked.
This prevents guest devices from accessing the protected server.

\subsection{Authorised Staff Access}

The authorised staff network is permitted to access the protected SSH
service on TCP port 22.

\subsection{Other Inbound Access}

Inbound access to the protected SSH service from sources outside the
authorised staff network is blocked.

The general policy is therefore:

\begin{center}
\begin{tabular}{lll}
\textbf{Source} & \textbf{Destination} & \textbf{Action} \\
\hline
Guest Network & Student Records Server & BLOCK \\
Staff Network & SSH Port 22 & ALLOW \\
Other Sources & SSH Port 22 & BLOCK \\
\end{tabular}
\end{center}

The firewall configuration is documented in the project file
\texttt{firewall.conf}.

\section{Firewall Test Procedure and Results}

Three test cases were defined to verify the firewall access-control
policy. One test checks permitted access and two tests check blocked
access.

\subsection{Test 1 -- Authorised Staff Connection}

The first test verifies that an authorised staff computer can connect to
the SSH service.

\begin{lstlisting}
nc -vz 192.168.10.10 22
\end{lstlisting}

\textbf{Expected outcome:} The connection should be permitted because
the source belongs to the authorised staff network.

\textbf{Simulated laboratory result:}

\begin{lstlisting}
Connection to 192.168.10.10 port 22 [tcp/ssh] succeeded!
\end{lstlisting}

The result indicates that the authorised connection was permitted.

\subsection{Test 2 -- Guest Network Connection}

The second test verifies that a guest network connection is blocked.

\begin{lstlisting}
nc -vz 192.168.10.10 22
\end{lstlisting}

\textbf{Expected outcome:} The connection should be blocked because the
source belongs to the guest network.

\textbf{Simulated laboratory result:}

\begin{lstlisting}
nc: connectx to 192.168.10.10 port 22 failed: Operation timed out
\end{lstlisting}

The result indicates that the guest connection was blocked.

\subsection{Test 3 -- Unauthorised Inbound Connection}

The third test verifies that another unauthorised source cannot access
the protected SSH service.

\begin{lstlisting}
nc -vz 192.168.10.10 22
\end{lstlisting}

\textbf{Expected outcome:} The connection should be blocked because the
source is outside the authorised staff network.

\textbf{Simulated laboratory result:}

\begin{lstlisting}
nc: connectx to 192.168.10.10 port 22 failed: Operation timed out
\end{lstlisting}

The result indicates that the unauthorised connection was blocked.

\subsection{Firewall Test Summary}

\begin{table}[h]
\centering
\begin{tabular}{llll}
\toprule
\textbf{Test} & \textbf{Source} & \textbf{Expected} & \textbf{Result} \\
\midrule
1 & Authorised Staff & ALLOW & PASS \\
2 & Guest Network & BLOCK & PASS \\
3 & Other Network & BLOCK & PASS \\
\bottomrule
\end{tabular}
\caption{Firewall test summary}
\end{table}

\textbf{Note:} The firewall results presented in this report are
simulated laboratory results because no specific network addresses or
service were supplied by the assessor.

\section{Conclusion}

The project demonstrated several important network security principles.

The risk assessment identified weaknesses involving weak passwords,
outdated software and unencrypted file transfers. Appropriate controls
were recommended, including strong passwords, multi-factor
authentication, regular software updates and encrypted file transfers.

The cryptography component demonstrated encryption and decryption for
protecting information. The integrity component demonstrated the use of
cryptographic hashing to detect changes to data.

The firewall component demonstrated network access control by allowing
authorised staff access to a protected service while blocking guest and
other unauthorised connections.

Together, these components demonstrate how confidentiality, integrity,
risk management and network access control can be applied as part of a
basic information security solution.

\section{GitHub Repository}

The complete project, including the source code, risk assessment,
firewall configuration, test evidence, README documentation and LaTeX
report, is maintained in the GitHub repository.

\textbf{Repository URL:}

\url{https://github.com/fabuwase/cryptography-network-security-exam}

\end{document}
```
